% Generated by Listas/make_statements.py from pratica-lecture-02.md, item 6. Do not edit.
\textbf{Source.} Muñoz-Garcia \& Toro-Gonzalez (2019), \emph{Strategy and Game Theory: Practice Exercises with Answers}, Ch. 3, Exercise 1 ``Game of Chicken'', pp.~86--87 {[}PDF p.~99--100{]}. Answer in the source: yes, a worked answer follows the exercise. Reproduced verbatim (statement only; the book's worked answer follows it).



\begin{quote}

Consider the game of Chicken (depicted in Fig. 3.1), in which two players driving their cars against each other must decide whether or not to swerve.



\textbf{Part (a)} Is there any strictly dominated strategy for Player 1? And for Player 2?



\textbf{Part (b)} What are the best responses for Player 1? And for Player 2?



\textbf{Part (c)} Can you find any pure strategy Nash Equilibrium (psNE) in this game?



\textbf{Part (d)} Find the mixed strategy Nash Equilibrium (msNE) of the game. \emph{Hint}: denote by \(p\) the probability that Player 1 chooses Straight and by \((1 - p)\) the probability that he chooses to Swerve. Similarly, let \(q\) denote the probability that Player 2 chooses Straight and \((1 - q)\) the probability that she chooses to Swerve.



\textbf{Part (e)} Show your result of part (d) by graphically representing every player \(i\)'s best response function \(BRF_i(s_j)\), where \(s_j = \{Swerve, Straight\}\) is the strategy selected by player \(j \neq i\).



{\def\LTcaptype{none} % do not increment counter

\begin{longtable}[]{@{}lll@{}}

\toprule\noalign{}

\emph{Player 1} ~\emph{Player 2} & Straight & Swerve \\

\midrule\noalign{}

\endhead

\bottomrule\noalign{}

\endlastfoot

Straight & 0,0 & 3,1 \\

Swerve & 1,3 & 2,2 \\

\end{longtable}

}



\emph{Fig. 3.1} Game of Chicken

\end{quote}

